\section*{Chapter \ref{ch:b3:exploration-seismology} --- Exploration Seismology}

\begin{solution}{exo:b3:exploration-seismology:1}
$Z = \rho v$: chalk $2000\times2500 = \qty{5.00e6}{kg.m^{-2}.s^{-1}}$, shale
$2300\times2300 = \qty{5.29e6}{kg.m^{-2}.s^{-1}}$, limestone $2600\times4500 =
\qty{11.70e6}{kg.m^{-2}.s^{-1}}$. Chalk--shale: $R = (5.29 - 5.00)/(5.29 + 5.00)
= 0.29/10.29 = 0.0282$; shale--limestone: $R = (11.70 - 5.29)/(11.70 + 5.29) =
6.41/16.99 = 0.377$. The reflected energy is $R^2 = 0.142$: the top of the
limestone sends back \qty{14.2}{\percent} of the incident energy, the chalk--shale
interface less than \qty{0.1}{\percent}.
\end{solution}

\begin{solution}{exo:b3:exploration-seismology:2}
$t_0 = 2\times800/2000 = \qty{0.800}{s}$; $t = \sqrt{0.800^2 + (1000/2000)^2} =
\sqrt{0.89} = \qty{0.9434}{s}$; the moveout is \qty{0.1434}{s}. The
approximation gives $1000^2/(2\times2000^2\times0.800) = \qty{0.1563}{s}$, too
large by \qty{9}{\percent}: at an offset larger than the depth, $x \ll vt_0$ no
longer holds.
\end{solution}

\begin{solution}{exo:b3:exploration-seismology:3}
Midpoints move by half the receiver interval: the CMP interval is
$12.5/2 = \qty{6.25}{m}$. The fold is $240\times12.5/(2\times25) = 60$. A fold of
120 needs a shot every $240\times12.5/(2\times120) = \qty{12.5}{m}$.
\end{solution}

\begin{solution}{exo:b3:exploration-seismology:4}
$\lambda = v/f$: \qty{100}{m} in the sandstone, $\lambda/4 = \qty{25}{m}$;
\qty{150}{m} in the limestone, $\lambda/4 = \qty{37.5}{m}$. With
$\lambda = 2800/30 = \qty{93.3}{m}$ and $z = \qty{1000}{m}$,
$r_F \approx \sqrt{\lambda z/2} = \sqrt{46\,667} = \qty{216}{m}$ (\qty{217}{m}
with the term $\lambda^2/16$).
\end{solution}

\begin{solution}{exo:b3:exploration-seismology:5}
Dix: $v^2 = (2800^2\times0.700 - 2400^2\times0.500)/0.200 = (5\,488\,000 -
2\,880\,000)/0.200 = \qty{1.304e7}{m^2.s^{-2}}$, so $v = \qty{3611}{m.s^{-1}}$.
The layer is crossed in \qty{0.100}{s} one way: $h = 3611\times0.100 =
\qty{361}{m}$.
\end{solution}

\begin{solution}{exo:b3:exploration-seismology:6}
(a) Let $d$ be the distance from the midpoint M to the reflector, so that
$t_0 = 2d/v$. The reflected ray from the source S to the receiver G is as long
as the straight line from the image S$'$ of S in the reflector to G. Put M at
the origin, the line along the dip; S$'$ lies at the distance $2d_S$ from S along
the normal, with $d_S = d - (x/2)\sin\varphi$. The square of S$'$G is
$x^2 + 4d_S^2 + 4d_Sx\sin\varphi = x^2 + 4d^2 - x^2\sin^2\varphi$, all the terms
in $dx$ cancelling, so $v^2t^2 = 4d^2 + x^2\cos^2\varphi$ and
$t^2 = t_0^2 + x^2\cos^2\varphi/v^2$. (b) The gather follows a hyperbola of speed
$V = v/\cos\varphi = 2000/\cos\ang{12} = \qty{2045}{m.s^{-1}}$: a dip makes the
\omterm{def:b3:exploration-seismology:velocities}{stacking velocity} too high, here by \qty{2.2}{\percent}.
\end{solution}

\begin{solution}{exo:b3:exploration-seismology:7}
By $\sqrt{60} = 7.75$. Doubling the signal-to-noise ratio needs four times as
many traces: a fold of 240.
\end{solution}

\begin{solution}{exo:b3:exploration-seismology:8}
The plotted depth $vt/2$ changes by $\qty{1.5}{km.s^{-1}}\times\qty{0.40}{s}
= \qty{0.60}{km}$ per kilometre: $\tan\theta_a = 0.60$, $\theta_a = \ang{31.0}$.
By the zero-offset theorem $\sin\theta = 0.60$: $\theta = \ang{36.9}$.
\end{solution}

\begin{solution}{exo:b3:exploration-seismology:9}
Each spike places a copy of $w$ centred on it: $0.3\,w$ on samples 0--2 gives
$(-0.15, 0.3, -0.15)$, and $-0.2\,w$ on samples 2--4 gives $(0.1, -0.2, 0.1)$.
Their sum is $(0, -0.15, 0.3, -0.05, -0.2, 0.1, 0)$, written from one sample
before the series to one after it. The reflections peak at samples 1 and 3, with
their true signs; they overlap on sample 2, where $-0.15 + 0.1 = -0.05$: a
little interference, but two distinct peaks.
\end{solution}

\begin{solution}{exo:b3:exploration-seismology:10}
$\arcsin(1800/2900) = \ang{38.37}$ and $\arcsin(1800/4400) = \ang{24.15}$, so
$i_c = (38.37 + 24.15)/2 = \ang{31.26}$ and $\theta = (38.37 - 24.15)/2 =
\ang{7.11}$; $v_2 = 1800/\sin\ang{31.26} = \qty{3469}{m.s^{-1}}$. The harmonic
mean $2\times2900\times4400/(2900 + 4400) = \qty{3496}{m.s^{-1}}$ is
\qty{0.8}{\percent} higher; multiplied by $\cos\theta = 0.9923$ it gives
\qty{3469}{m.s^{-1}}, the usual small-dip approximation
$v_2 \approx \cos\theta\,2v_dv_u/(v_d + v_u)$.
\end{solution}

\begin{solution}{exo:b3:exploration-seismology:11}
(a) The composite is $R[w(t) - w(t - \Delta t)] \approx R\,\Delta t\,w'(t)$ by
the first-order Taylor expansion. (b) Its amplitude is proportional to
$\Delta t$, so a bed thinner than the \omterm{def:b3:exploration-seismology:resolution}{tuning thickness} still returns a signal,
with the shape of $w'$, but its top and base cannot be picked as two events: it
is detected, not resolved. If $R$ and the wavelet are known, for instance from a
well, the amplitude gives $\Delta t$ and so the thickness. (c) The amplitude is
proportional to the thickness: the ratio is $5/10 = 1/2$.
\end{solution}

\begin{solution}{exo:b3:exploration-seismology:12}
(a) $t(x) = (2/v)\sqrt{z^2 + x^2}$: $t(0) = 2\times600/3000 = \qty{0.400}{s}$ and
$t(300) = (2/3000)\sqrt{600^2 + 300^2} = \qty{0.4472}{s}$. (b) $t^2 = t_0^2 +
x^2/(v/2)^2$, the NMO form with the speed $v/2$. A reflection with the same apex
time, on a zero-offset section, would show no moveout at all; compared with an
NMO hyperbola of the same $t_0$ in a gather, the diffraction's moveout
$x^2/(2(v/2)^2t_0)$ is four times larger, because the time depends on the
distance from the scatterer twice, down and back. (c) Migration sums the
section along this hyperbola and puts the energy at its apex: the diffraction
collapses to a point at $x = 0$, $z = \qty{600}{m}$.
\end{solution}

\begin{solution}{pb:b3:exploration-seismology:1}
\textbf{1.} From W, $400/(486.1 - 425.2) = 400/0.0609 = \qty{6568}{m.s^{-1}}$;
from E, $400/(486.1 - 372.4) = 400/0.1137 = \qty{3518}{m.s^{-1}}$. The head
wave from W seems fast: W is the down-dip end, and its shots travel up-dip; the
refractor rises from W to E.

\textbf{2.} $\arcsin(2400/3518) = \ang{43.02}$, $\arcsin(2400/6568) =
\ang{21.43}$: $i_c = \ang{32.22}$, $\theta = \ang{10.79}$.

\textbf{3.} $v_2 = 2400/\sin\ang{32.22} = \qty{4501}{m.s^{-1}}$: within the
range of limestone, the top of the Jurassic limestone.

\textbf{4.} The intercepts are $t - x/v_{\text{app}}$ at \qty{1200}{m}:
$t_{i,W} = 0.4861 - 1200/6568 = \qty{0.3034}{s}$ and $t_{i,E} = 0.4861 -
1200/3518 = \qty{0.1450}{s}$. With $\cos i_c = 0.8460$, $h = v_1t_i/(2\cos
i_c)$: $h_W = 2400\times0.3034/1.692 = \qty{430.4}{m}$ and $h_E =
\qty{205.7}{m}$.

\textbf{5.} Dividing by $\cos\theta = 0.9823$: \qty{438.1}{m} under W and
\qty{209.4}{m} under E. Then $(438.1 - 209.4)/1200 = 0.1906 =
\tan\ang{10.79}$, the dip of question 2.

\textbf{6.} The shale returns no head wave: it is invisible, and the survey
counts it with the chalk. Here $v_1 = \qty{2.40}{km.s^{-1}}$ was given as the
mean of the two; had the chalk's direct-wave speed, \qty{2.50}{km.s^{-1}},
been used for the whole overburden, every depth would come out too large.

\textbf{7.} $V = 600/\sqrt{0.5437^2 - 0.4878^2} = 600/0.2401 =
\qty{2499}{m.s^{-1}}$ and $V = 600/\sqrt{0.6240^2 - 0.5748^2} =
\qty{2471}{m.s^{-1}}$. The second averages the chalk with the slower shale
below it: the mean drops.

\textbf{8.} \qty{2811}{m.s^{-1}} and \qty{2829}{m.s^{-1}}.

\textbf{9.} Layer 1: \qty{2499}{m.s^{-1}}. Layer 2: $v^2 = (2471^2\times0.5748
- 2499^2\times0.4878)/0.0870$, $v = \qty{2308}{m.s^{-1}}$. Layer 3:
$v^2 = (2811^2\times0.6548 - 2471^2\times0.5748)/0.0800$,
$v = \qty{4561}{m.s^{-1}}$. Layer 4: $v^2 = (2829^2\times0.7348 -
2811^2\times0.6548)/0.0800$, $v = \qty{2972}{m.s^{-1}}$.

\textbf{10.} Chalk (\qty{2499}{m.s^{-1}}), shale (\qty{2308}{m.s^{-1}}),
limestone (\qty{4561}{m.s^{-1}}), sandstone (\qty{2972}{m.s^{-1}}), each within
its published range. The reflections are the chalk--shale interface, the top
of the limestone, the limestone--sandstone interface and the unconformity at
the top of the schist basement.

\textbf{11.} Chalk and shale have almost the same impedance: $R = 0.028$
(exercise 1), against 0.38 for the top of the limestone.

\textbf{12.} The hyperbola of speed $V_{\text{rms}}$ holds only to second order
in the offset; at offsets of \qty{2}{km}, twice the depth of the deepest
reflector, the times depart from it and a hyperbola fitted to the whole spread
would give wrong velocities, and Dix intervals still worse. The far offsets
also stretch the wavelets when corrected.

\textbf{13.} Two-way times in the layers: \qty{0.4878}{s}, \qty{0.0870}{s},
\qty{0.0800}{s}, \qty{0.0800}{s}. Thicknesses $v\,\Delta t/2$: \qty{609.5}{m},
\qty{100.4}{m}, \qty{182.4}{m}, \qty{118.9}{m}.

\textbf{14.} Depths below the ground: \qty{609.5}{m}, \qty{709.9}{m},
\qty{892.3}{m}, \qty{1011.2}{m}. The shale, limestone and sandstone come out
within \qty{0.4}{m}, \qty{2.4}{m} and \qty{1.1}{m} of their thicknesses, the
error growing in the fast, thin limestone where a small error of time costs
most.

\textbf{15.} $\bar v = 2\times1011.2/0.7348 = \qty{2752}{m.s^{-1}}$. The RMS
velocity weights each layer by $v^2$ and is never less than the time-weighted
mean: $2829\times0.7348/2 = \qty{1039}{m}$, \qty{28}{m} too deep.

\textbf{16.} $\lambda = 2308/30 = \qty{76.9}{m}$, $\lambda/4 = \qty{19.2}{m}$;
the Ricker estimate is $2308/(5.2\times30) = \qty{14.8}{m}$. The shale is
\qty{100}{m} thick, five times its \omterm{def:b3:exploration-seismology:resolution}{tuning thickness}: its top and base are
resolved. It would tune, by the quarter-wavelength rule, at $f = v/(4b) =
2308/400 = \qty{5.77}{Hz}$.

\textbf{17.} $\lambda = 2829/30 = \qty{94.3}{m}$, $z = \qty{1011}{m}$:
$r_F \approx \sqrt{94.3\times1011/2} = \qty{218}{m}$ (\qty{220}{m} with the
term $\lambda^2/16$).

\textbf{18.} $t(x)^2 = t_0^2 + (2x/v)^2$, so $v = 2x/\sqrt{t^2 - t_0^2} =
400/\sqrt{0.5732^2 - 0.550^2} = 400/0.1614 = \qty{2478}{m.s^{-1}}$.

\textbf{19.} $z = vt_0/2 = 2478\times0.550/2 = \qty{681}{m}$.

\textbf{20.} The plotted depth $vt/2$ changes by $1.24\times0.167 = 0.207$ per
unit length: $\theta_a = \arctan0.207 = \ang{11.70}$. The true dip has
$\sin\theta = 0.207$: $\theta = \ang{11.95}$, the dip of the limb.

\textbf{21.} $d = 2480\times0.655/2 = \qty{812}{m}$; migration moves the
reflection $d\sin\theta = 812\times0.207 = \qty{168}{m}$ up-dip, towards the
crest of the anticline, and up from \qty{812}{m} to $d\cos\theta =
\qty{795}{m}$.

\textbf{22.} $2500\times0.160/2 = \qty{200}{m}$.

\textbf{23.} From questions 13 and 14, the basement unconformity lies
\textbf{\qty[text-series-to-math=true]{1011}{m}} below the ground at CMP S.
\end{solution}
